Vairumā temporālās kodēšanas gadījumu pēdējā slānī izmanto impulsu akumulatoru \cite{Fang2021, kim2021, rathi2020}. Galīgo izeju aprēķina kā vairāku impulsu summu visos \(T\) laika soļos, lai veicinātu visaugstāko impulsu ģenerēšanas biežumu neironam, kas saistīts ar sagaidāmo izeju.

Impulsu pamattīkla izejas impulsi tiek pārveidoti uzkrātā pazīmju reprezentācijā, pirms tos nodod \Gls{ssd} detektēšanas galvai. Izeja tiek vidējota visos \(T\) laika soļos. Izejas pazīmju kartei ar kanālu \(c\) un telpiskajām koordinātām \(h, w\) laikā vidējoto pazīmju reprezentāciju nosaka
\begin{equation}
    F_{c,h,w} =
    \frac{1}{T}\sum_{t=0}^{T-1} S_{t,c,h,w},
\end{equation}
kur \(S_{t,c,h,w}\in\{0,1\}\) ir neirona izejas impulss laika solī \(t\). Iegūtā \(F\) reprezentē impulsu pazīmju kartes vidējo impulsu aktivitāti. Pēc tam detektēšanas galva veic objektu klasifikāciju, ietverošo taisnstūru regresiju un orientācijas novērtēšanu.

Ietverošo taisnstūru detektēšanai un klasifikācijai izmanto \Gls{ssd} \cite{liu2016} prognozēšanas galvas. Tās saņem \Gls{snn} pamattīkla ģenerētās pazīmju kartes, nodrošinot vienlaicīgu objektu detektēšanu un to ietverošo taisnstūru regresiju. Paplašinot iepriekšējos darbus \cite{Cordone2022, fan2025}, arhitektūra papildināta ar orientācijas regresiju, kas piedāvāta darbā \cite{simony2018}, lai detektētu vairāku klašu orientētus ietverošos taisnstūrus. Katram detektētajam objektam detektēšanas galva ģenerē izejas vektoru
\begin{equation}
    d = [p_c, b_x, b_y, b_w, b_l, b_\theta, c],
\end{equation}
kur \(p_c\) ir prognozētā ticamības vērtība, \(b_x\) un \(b_y\) ir ietverošā taisnstūra centra koordinātas \Gls{bev} plaknē, \(b_w\) un \(b_l\) ir tā platums un garums, \(b_\theta\) ir tā orientācijas leņķis attiecībā pret \Gls{bev} koordinātām un \(c\) ir prognozētā objekta klase.

Tīkls tiek apmācīts ar daudzuzdevumu mērķfunkciju, kas apvieno klasifikācijas, ietverošā taisnstūra regresijas un orientācijas regresijas zudumus:
\begin{equation}
  \mathcal{L} =
  \mathcal{L}_{cls}
  + \lambda_{box}\mathcal{L}_{box}
  + \lambda_{ang}\mathcal{L}_{ang},
\end{equation}
kur \(\mathcal{L}_{cls}\) ir objektu klasifikācijas zudums, \(\mathcal{L}_{box}\) ir ietverošā taisnstūra regresijas zudums un \(\mathcal{L}_{ang}\) ir orientācijas regresijas zudums. Koeficienti \(\lambda_{box}\) un \(\lambda_{ang}\) nosaka regresijas locekļu ieguldījumu kopējā apmācības mērķfunkcijā.

Klasifikācijas zudumam \(\mathcal{L}_{cls}\) izmanto fokālo zuduma funkciju \cite{lin2017}, lai samazinātu vieglāk klasificējamu negatīvo piemēru ieguldījumu un risinātu klašu nelīdzsvarotību. Ietverošā taisnstūra un orientācijas regresijai izmanto \textit{Smooth L1} zuduma funkciju tās noturības pret izlecošām vērtībām un stabilās optimizācijas uzvedības dēļ. Tādējādi apvienotā mērķfunkcija apmāca impulsu pamattīklu un detektēšanas galvas iegūt pazīmes un vienlaikus optimizēt objektu lokalizāciju un orientāciju.

\Gls{ssd} klasifikācijas un regresijas galvas veido viens konvolūcijas slāņu posms bez papildu aktivācijas funkcijām starp
tiem. Tas ir īpaši piemēroti ar aparatūru paātrinātiem \Gls{snn}, jo netiek pievienoti aktivācijas posmi, kuriem būtu
nepieciešama papildu stāvokļu glabāšana un aprēķini, tādējādi nodrošinot saderību ar uz notikumiem balstītu apstrādi \Gls{fpga}
realizācijās. Šādās realizācijās svarīgi, lai aparatūras paātrinātājs atbalstītu papildu savienojumus, kas nodrošina piekļuvi
dažādos pamattīkla līmeņos iegūtām daudzlīmeņu pazīmju kartēm.

Lai saglabātu objektu identitāti starp kadriem, detektors apvienots ar vairāku objektu izsekošanas algoritmu
\textit{AB3DMOT} \cite{weng2020}. Izsekotājs no detektēšanas galvas saņem prognozētos orientētos ietverošos taisnstūrus,
klašu marķējumus un ticamības vērtības un sasaista detektējumus starp kadriem, izmantojot konfigurējamu telpiskās līdzības
metriku. Trajektorija tiek apstiprināta pēc noteikta veiksmīgu sasaistījumu skaita un saglabāta ierobežotu kadru skaitu bez
atbilstoša detektējuma.


